\section{Boundaries}

The plan boundary is the governing claim: Bayesian outputs are analytic unless
explicitly calibrated to a risk limit. The implementation follows that line. It
verifies a declared Bayesian audit record, but it does not perform a full
posterior computation and does not certify that the posterior upset probability
has frequentist risk-limiting coverage.

Prior dependence is part of the public contract. The same tabulation evidence
can produce different posterior upset probabilities under different priors or
likelihood models, so the verifier requires non-empty identifiers for both. It
still treats those identifiers as declared metadata. A future implementation
could bind them to executable posterior code, hand-computable fixtures, or an
external calibration proof; V.20 does not yet do that.

Bayesian audits are therefore not risk-limiting in the frequentist sense merely
because they report a small upset probability. RCOUNT keeps the field
\texttt{calibrated\_risk\_limit\_ppm} optional and shape-checked only. Until a
calibration rule is added, that field is a boundary marker, not a certificate.
This is also why the replay transcript reports Bayesian tabulation as recorded
posterior analytics rather than as BRAVO, Kaplan-Markov, ALPHA, or another RLA
stopping calculation.
